This paper has built the predictive distribution of a passenger ship's mission
failure count from a hierarchical model with crossed ship-wide and ship--product
terms, fitted from pooled or divided records, and has followed it through to the
spares a ship embarks.

For one population, the negative binomial matched to the moments of the mixture
is close to the exact plan, at most \spNbShort{} points short at the 99\,\%
level on the grid examined, and can be used where the exact computation is not
available. The Poisson plan and the variance-inflated margin should not be. For
a zone with several references, the plan is a vector whose cover depends on how
the heterogeneity divides between ship and product. Per-reference quantiles fall
short of a joint target, cost-minimal vectors save a median \rvSaveMed\,\% over
a common stock, and while the division is unknown a conservative plan keeps the
target at a premium in stock.

In the simulated passenger fleet the largest loss came from planning with
datasheet rates, which covered \flMtbfJointLo--\flMtbfJointHi\,\% of missions
against a 99\,\% target even as a joint plan. Joint plans fitted to the fleet's
own records covered \flStatLo--\flStatHi\,\% under supply events, bad lots,
heavy-tailed effects and misrecorded replacements, and from half the fleet over
one year. With the moderate heterogeneity assumed there, the hierarchy added
little to a pooled negative binomial on average, but it kept its stock in
proportion when one ship lay far in the tail, and the pooled model did not. On
the street-lighting network, where the
interaction between area and family was more than twice as large, the
hierarchical plans were the only ones to meet the target in the held-out year,
with fewer units than the pooled negative binomial in known areas and at a
high price for an area never seen.

What the model needs from a fleet's records is set out in
Section~\ref{sec:needed}: several references followed on the same ships over
repeated missions, ten ships or more, exposure in powered hours and a record of
common-cause events. Applying it to such records is the next step.
